
% ---------- Méthodes existantes : RO classique ----------
\begin{frame}{Méthodes existantes : recherche opérationnelle}
  \textbf{Trois niveaux de décision, largement étudiés séparément :}
  \vspace{0.3em}
  \begin{table}
    \centering\small
    \begin{tabular}{lll}
      \toprule
      \textbf{Niveau} & \textbf{Horizon} & \textbf{Décisions} \\
      \midrule
      Stratégique & 1--5 ans & Dimensionnement des salles \\
      Tactique & 1--3 mois & Programme maître de chirurgie (MSS) \\
      Opérationnel & jour -- semaines & Affectation salle / patient / lit \\
      \bottomrule
    \end{tabular}
  \end{table}
  \begin{block}{Formulation dominante}
    \textbf{Programmation linéaire (mixte)} : variables d'affectation
    patient~$\rightarrow$~bloc~$\rightarrow$~lit, sous contraintes de capacité
    et de durées, fonction objectif minimisant retards et dépassements.
  \end{block}
\end{frame}

% ---------- Heuristiques vs métaheuristiques ----------
\begin{frame}{Heuristiques et métaheuristiques}
  \small
  Le problème d'ordonnancement bloc/lits est \alert{NP-difficile} : la
  programmation linéaire exacte devient vite trop coûteuse en calcul. Deux
  familles de méthodes approchées coexistent.
  \vspace{0.2em}
  \begin{columns}[T]
    \begin{column}{0.48\textwidth}
      \begin{block}{Heuristique}
        \small
        \begin{itemize}
          \item Règle empirique conçue pour \textbf{un problème précis}
          \item Rapide, peu coûteuse en calcul
          \item Aucune garantie d'optimalité : risque d'optimum local
        \end{itemize}
      \end{block}
    \end{column}
    \begin{column}{0.48\textwidth}
      \begin{block}{Métaheuristique}
        \small
        \begin{itemize}
          \item Cadre \textbf{générique}, applicable à des problèmes
          NP-difficiles variés
          \item Explore plus largement l'espace des solutions
          \item Meilleure qualité de solution, au prix d'un calcul plus long
        \end{itemize}
      \end{block}
    \end{column}
  \end{columns}
\end{frame}

% ---------- Limites + synthèse ----------
\begin{frame}{Limites identifiées dans la littérature}
  \begin{alertblock}{Limite principale}
    Blocs opératoires et lits d'hospitalisation sont \textbf{le plus souvent
    optimisés séparément} : le bloc est planifié, puis le lit est recherché
    a posteriori.
  \end{alertblock}
  \begin{itemize}
    \item Peu de modèles \textbf{conjoints} bloc + lit dans les approches
    classiques (programmation linéaire).
    \item Les incertitudes (urgences, variabilité des durées) sont rarement
    intégrées de façon robuste dans les modèles conjoints.
  \end{itemize}
  \begin{center}
    \large \textcolor{mcmPrimary}{\textbf{$\Rightarrow$ Optimisation conjointe blocs + lits}}
  \end{center}
\end{frame}
